% ============================================================
% Propuesta de PPS - "Simulador Concurrente"
% Federico Dobal - Analista en TIC, Facultad de Informatica UNLP
%
% Como compilar: cualquier distribucion TeX (pdflatex) o subir el
% .tex directo a Overleaf, ahi compila sin instalar nada.
%
% Las capturas van en recorrido_visual.tex: reemplazar los
% recuadros grises "PLACEHOLDER" por \includegraphics, poniendo
% los archivos de imagen en esta misma carpeta.
% ============================================================
\documentclass[11pt,a4paper]{article}

\usepackage[utf8]{inputenc}
\usepackage[T1]{fontenc}
\usepackage[spanish,es-noshorthands]{babel}
\usepackage{lmodern}
\usepackage[margin=2.7cm]{geometry}
\usepackage{parskip}
\usepackage{xcolor}
\usepackage{tikz}
\usetikzlibrary{arrows.meta, positioning, shapes.rounded}
\usepackage{enumitem}
\usepackage{graphicx}
\usepackage{tcolorbox}
\usepackage{titlesec}
\usepackage{hyperref}

\definecolor{unlpblue}{HTML}{1B3B6F}
\definecolor{boxgray}{HTML}{F4F5F7}
\definecolor{boxline}{HTML}{C9CDD3}

\hypersetup{
    colorlinks=true,
    linkcolor=unlpblue,
    urlcolor=unlpblue,
    pdftitle={Propuesta PPS - Simulador Concurrente},
    pdfauthor={Federico Dobal}
}

\titleformat{\section}{\large\bfseries\color{unlpblue}}{\thesection.}{0.5em}{}
\titleformat{\subsection}{\normalsize\bfseries\color{unlpblue}}{\thesubsection.}{0.5em}{}
\titlespacing*{\section}{0pt}{1.4em}{0.7em}
\titlespacing*{\subsection}{0pt}{1.1em}{0.5em}

\setlength{\parindent}{0pt}
\setlength{\parskip}{0.6em}

\begin{document}

% ---------- Portada / encabezado ----------
\begin{center}
    {\Large\bfseries\color{unlpblue} Simulador Concurrente}\\[0.3em]
    {\large Propuesta de Práctica Profesional Supervisada}\\[0.6em]
    {\normalsize Carrera Analista en TIC --- Facultad de Informática, UNLP}
\end{center}

\vspace{0.8em}

\begin{tabular}{ll}
    \textbf{Alumno:} & Federico Dobal \\
    \textbf{Destinatario:} & Prof. Franco Chichizola \\
    \textbf{Fecha:} & 28 de julio de 2026 \\
    \textbf{Demo en vivo:} & \url{https://simulador-concurrente.vercel.app/}
\end{tabular}

\vspace{1em}
\hrule
\vspace{1em}

% ---------- Introducción ----------
\section{Introducción}

Este documento acompaña una demostración en vivo de \textbf{Simulador Concurrente}, una
herramienta pensada para escribir y ejecutar programas concurrentes directamente desde el
navegador. La idea es proponerlo como tema de Práctica Profesional Supervisada: ya existe
un desarrollo funcional (el que se muestra en la demo), y lo que se busca es continuar su
evolución --- nuevas construcciones del lenguaje, mejoras en la visualización, casos de
prueba guiados --- bajo la supervisión de la cátedra.

% ---------- Motivación ----------
\section{Motivación}

Durante la cursada de Programación Concurrente es habitual notar la falta de herramientas
para probar y verificar la correctitud de un programa concurrente antes de entregarlo. Los
ejercicios se resuelven en un pseudocódigo que, aunque sigue las convenciones de la
cátedra, termina siendo bastante subjetivo a quien lo escribe: no hay forma de correrlo y
ver si realmente funciona.

A eso se suma que razonar sobre exclusión mutua, semáforos o barreras exige imaginar
simultáneamente varios hilos avanzando en un orden que no se puede predecir. Es un
ejercicio mental que cuesta hacer únicamente sobre el papel, y donde una condición de
carrera puede pasar completamente desapercibida hasta que alguien la señala.

Simulador Concurrente busca achicar esa brecha: se escribe el pseudocódigo tal como se
define en la cátedra, se ejecuta paso a paso y se observa en tiempo real cómo evoluciona
cada hilo, cada variable compartida y cada estructura de datos.

% ---------- Descripción ----------
\section{Descripción del simulador}

Es una aplicación web de una sola página, construida en React 19 y TypeScript sobre Vite.
No depende de ningún servidor: todo el análisis y la ejecución del pseudocódigo ocurren en
el propio navegador, lo que la hace liviana y fácil de usar en una clase o un TP sin
instalar nada.

El lenguaje soportado cubre lo que normalmente aparece en los ejercicios de la materia:
variables (\texttt{int}, \texttt{boolean}, \texttt{double}, \texttt{text}), semáforos y
arreglos de semáforos, colas FIFO, procesos parametrizados (\texttt{process P[i: 0..4]}
para instanciar varios hilos de una), bloques atómicos, espera condicional
(\texttt{await}, la base para monitores) y las estructuras de control usuales
(\texttt{if}, \texttt{while}, \texttt{for}).

\begin{tcolorbox}[colback=boxgray, colframe=boxline, boxrule=0.5pt, arc=2pt, left=8pt, right=8pt, top=6pt, bottom=6pt]
\texttt{\footnotesize
sem mutex = 1;\\
int comensales = 0;\\[2pt]
process Filosofo[i: 0..4] \{\\
\phantom{xx}< await(comensales < 4) comensales++; >\\
\phantom{xx}P(mutex);\\
\phantom{xx}\ldots\\
\phantom{xx}V(mutex);\\
\}
}
\end{tcolorbox}

\subsection{Arquitectura del proyecto}

El código se organiza en dos capas bien separadas: un \textbf{motor de simulación} que no
sabe nada de interfaces gráficas (podría reutilizarse tal cual en otro contexto), y una
\textbf{interfaz} en React que se limita a mostrar lo que el motor va calculando.

\subsubsection*{Motor de simulación (\texttt{src/engine})}

Al presionar ``Ejecutar'', el código pasa por tres etapas antes de empezar a correr:

\begin{center}
\begin{tikzpicture}[
    node distance=0.9cm and 1.1cm,
    box/.style={draw=unlpblue, thick, rounded corners=3pt, fill=boxgray,
                minimum height=1.1cm, minimum width=2.5cm, align=center, font=\small},
    lbl/.style={font=\scriptsize\itshape, color=unlpblue!70!black}
]
    \node[box] (src)  {Código\\fuente};
    \node[box, right=of src]  (lex)  {Lexer};
    \node[box, right=of lex]  (par)  {Parser};
    \node[box, right=of par]  (vm)   {Máquina\\virtual};

    \draw[-{Latex[length=2mm]}, unlpblue] (src) -- (lex) node[midway, above, lbl] {tokenize()};
    \draw[-{Latex[length=2mm]}, unlpblue] (lex) -- (par) node[midway, above, lbl] {Token[]};
    \draw[-{Latex[length=2mm]}, unlpblue] (par) -- (vm)  node[midway, above, lbl] {AST};
\end{tikzpicture}
\end{center}

El \emph{lexer} convierte el texto en una lista de tokens (números, identificadores,
palabras clave, operadores), descartando espacios y comentarios. El \emph{parser} arma con
eso un árbol de sintaxis tipado, distinguiendo declaraciones globales, sentencias dentro de
un proceso y definiciones de proceso. Y la \emph{máquina virtual} interpreta ese árbol
directamente, sin compilar a ningún otro lenguaje intermedio.

Lo más interesante pasa ahí: en cada paso de la simulación, la máquina virtual baraja los
hilos que están listos para correr y elige uno al azar (incluso dejando que algún hilo
``pierda el turno'' de vez en cuando), para que la ejecución se sienta tan no determinística
como la de hilos reales. Los semáforos bloquean y despiertan hilos en orden FIFO, los
bloques atómicos se ejecutan de un tirón, y cada paso relevante queda registrado con una
foto completa del estado --- lo que permite después recorrer la ejecución hacia adelante y
hacia atrás sin tener que volver a correr nada.

\subsubsection*{Interfaz de usuario (\texttt{src/components})}

Del lado visual la cosa es más simple: un editor de código con indentación automática y
deshacer/rehacer, un panel que muestra el estado en vivo de cada hilo, variable y semáforo,
y una línea de tiempo estilo diagrama de Gantt donde se ve qué hizo cada proceso a lo largo
de la ejecución. Nada más --- la gracia está en que todo se actualiza solo, sin que el
usuario tenga que ir refrescando nada a mano.

\subsection{Recorrido visual}

% ------------------------------------------------------------
% Reemplazar cada recuadro por \includegraphics[width=\linewidth]{archivo.png}
% ------------------------------------------------------------

\textbf{Editor de código}

\begin{center}
\fbox{\parbox[c][5cm][c]{0.9\linewidth}{\centering\color{gray}\itshape
CAPTURA --- editor con el pseudocódigo cargado
}}
\end{center}

\textbf{Panel de estado en vivo}

\begin{center}
\fbox{\parbox[c][5cm][c]{0.9\linewidth}{\centering\color{gray}\itshape
CAPTURA --- hilos, variables y semáforos durante la ejecución
}}
\end{center}

\textbf{Línea de tiempo}

\begin{center}
\fbox{\parbox[c][5cm][c]{0.9\linewidth}{\centering\color{gray}\itshape
CAPTURA --- diagrama de Gantt de la ejecución
}}
\end{center}

\vspace{1.5em}
\hrule
\vspace{0.8em}

\begin{center}
\small La demo completa está disponible en \url{https://simulador-concurrente.vercel.app/}
\end{center}

\end{document}
